% Part: first-order-logic
% Chapter: syntax-and-semantics
% Section: Structures

\documentclass[../../../include/open-logic-section]{subfiles}

\begin{document}

\olfileid{fol}{syn}{str}

\olsection{\printtoken{P}{structure} للغات الرتبة الأولى}

\begin{explain}
لغة الرتبة الأولى إذا نُظر إليها في ذاتها كانت \emph{غير مؤوَّلة}؛
إذ لا معنى معيّنًا لرموز ال!!{constant}s وال!!{function}s
وال!!{predicate}s بمجردها. وإنما نعيّن معانيها باختيار
\article{structure} \emph{!!{structure}}. ففيها يتعيّن
\emph{المجال}: الأشياء التي تدل عليها ال!!{constant}s، وتجري
عليها ال!!{function}s، وتتراوح عليها المكممات. ويتعيّن فيها كذلك
مدلول كل واحد من ال!!{constant}s، وكيف تنقل !!a{function} الأشياء
إلى أشياء، وعلى أي الأشياء تصدق ال!!{predicate}s.
وعلى ال!!^{structure}s تُبنى المفاهيم \emph{الدلالية} في المنطق،
كالنتيجة والصلاحية وقابلية الإرضاء. وقد تختلف تسميتها في
المؤلفات، فيقال «بنى» أو «تأويلات» أو «نماذج».
\end{explain}

\begin{defn}[!!^{structure}s]
تتألف \Article{structure} \emph{!!{structure}}~$\Struct M$ للغة
$\Lang{L}$ من منطق الرتبة الأولى من العناصر الآتية:
\begin{enumerate}
\item \emph{المجال:} مجموعة غير خالية، $\Domain M$
\item \emph{تأويل ال!!{constant}s:} لكل !!{constant}~$c$ من
  $\Lang{L}$، !!a{element}~$\Assign{c}{M} \in \Domain M$
\item \emph{تأويل ال!!{predicate}s:} لكل !!{predicate} ذات
  $n$ مواضع~$R$ من~$\Lang{L}$ (عدا~$\eq$)، علاقة ذات $n$ مواضع
  $\Assign{R}{M} \subseteq \Domain{M}^n$
\item \emph{تأويل ال!!{function}s:} لكل !!{function} ذات
  $n$ مواضع~$f$ من~$\Lang{L}$، دالة ذات $n$ مواضع
  $\Assign{f}{M} \colon \Domain{M}^n \to \Domain{M}$
\end{enumerate}
\end{defn}

\begin{ex}
تتألف !!^a{structure}~$\Struct M$ للغة الحساب من مجموعة، ومن عنصر
من~$\Domain M$، هو $\Assign{\Obj 0}{M}$، تأويلًا لل!!{constant}
~$\Obj 0$، ومن دالة أحادية المحل
$\Assign{\Obj \prime}{M} \colon \Domain{M} \to \Domain M$،
ودالتين ثنائيتي المحل $\Assign{\Obj +}{M}$ و$\Assign{\Obj
  \times}{M}$، كلتاهما $\Domain M^2 \to \Domain M$، وعلاقة ثنائية
المحل $\Assign{\Obj <}{M} \subseteq \Domain{M}^2$.

وأوضح مثال لهذه البنية أن نختار ما يأتي:
\begin{enumerate}
\item $\Domain N = \Nat$
\item $\Assign{\Obj 0}{N} = 0$
\item $\Assign{\Obj \prime}{N}(n) = n + 1$ لكل $n \in \Nat$
\item $\Assign{\Obj +}{N}(n, m) = n + m$ لكل $n, m \in \Nat$
\item $\Assign{\Obj \times}{N}(n, m) = n\cdot m$ لكل $n, m \in \Nat$
\item $\Assign{\Obj <}{N} = \Setabs{\tuple{n, m}}{n \in \Nat, m \in
  \Nat, n < m}$
\end{enumerate}
وهذه البنية~$\Struct N$ للغة~$\Lang L_A$ تسمى
\emph{النموذج القياسي للحساب}؛ فإن الرموز غير المنطقية
لـ~$\Lang L_A$ تأخذ فيها المعاني الحسابية المعهودة.

وليست هذه كل ال!!{structure}s الممكنة لـ~$\Lang L_A$، بل غيرها كثير.
فلنا أن نجعل الأعداد الصحيحة~$\Int$ مجالًا مكان~$\Nat$، ثم نختار
تأويلات~$\Obj 0$، و$\Obj \prime$، و$\Obj +$، و$\Obj
\times$، و$\Obj <$ على وفق ذلك. ولنا أيضًا أن نضع
لـ~$\Lang L_A$ بنى لا صلة لها بالأعداد أصلًا.
\end{ex}

\begin{ex}
أما البنية~$\Struct M$ للغة نظرية المجموعات~$\Lang L_Z$، فيكفي
لتعيينها مجموعة غير خالية وعلاقة ثنائية المحل. فيجوز، بمقتضى
التعريف الصوري، أن نأخذ مجموعة الناس مع العلاقة «$x$ أكبر سنًا من~$y$»
بوصفهما !!a{structure} لـ~$\Lang L_Z$، كما يمكن استخدام~$\Nat$ مع
$n \ge m$، حيث $n, m \in \Nat$.

ومن ال!!{structure}s الجديرة بالنظر للغة~$\Lang L_Z$ ما تكون
!!{element}s مجاله مجموعات بالفعل، ويكون مدلول~$\Obj \in$
فيه نفس العلاقة «$x$ !!a{element} من~$y$»، ال!!{structure}
~$\Struct{{HF}}$ المؤلفة من \emph{المجموعات المنتهية وراثيًا}:
\begin{enumerate}
\item $\Domain{{{HF}}} = \emptyset \cup \Pow{\emptyset} \cup
  \Pow{\Pow{\emptyset}} \cup \Pow{\Pow{\Pow{\emptyset}}} \cup \dots$;
\item $\Assign{\Obj \in}{{{HF}}} = \Setabs{\tuple{x, y}}{x, y \in
  \Domain{{{HF}}}, x \in y}$.
\end{enumerate}
\end{ex}

\begin{digress}
وللشروط التي نضعها في تعريف !!a{structure} أثر في المنطق نفسه.
فمنع المجال الخالي يقتضي، مع المعنى المعتاد لإرضاء الجمل المكممة
أو صدقها، أن تكون
$\lexists[x][(!A(x) \lor \lnot !A(x))]$ صالحة، أي حقيقة منطقية.
واشتراط أن تدل ال!!{constant}s كلها على أشياء من المجال يجعل
التعميم الوجودي استدلالًا سليمًا: من $!A(a)$ يلزم
$\lexists[x][!A(x)]$. أما إن أجزنا أن تدل الأسماء على ما خارج المجال،
أو ألا تدل على شيء، فقد أخذنا بمسلك \emph{المنطق الحر}؛
وفيه يحتاج التعميم الوجودي إلى مقدمة زائدة: من $!A(a)$
و$\lexists[x][\eq[x][a]]$ يلزم
$\lexists[x][!A(x)]$.
\end{digress}

\end{document}
